% Related lecture: SC4002-L07
% Source: Week 7-Dependency Parsing, pp. 3, 6, 13--14, 17
% Lecture figures and worked example restated with source-checked explanations.

\exercisegroup{SC4002-L07-EX}{第 07 讲：Dependency Parsing}

\exercisemeta
  {SC4002-L07-EX}
  {2026-09-23}
  {Week 7-Dependency Parsing pp.~3, 6, 13--14, 17；本地字幕 69612}
  {E01--E03 为讲义例图与算法例题的重述及解析，非 Tutorial}
  {讲解后请求整理；图中关系与操作顺序已核对}

\exercisequestion{SC4002-L07-E01}{读依存树：航班修饰与双宾语}

\textbf{来源与题意：}讲义 pp.~3, 6 的例图。分析 \texttt{I prefer the morning flight through Denver} 与 \texttt{John gives Marry an apple}，区分 head、dependent 及关系类型；人名沿用讲义拼写。

\textbf{对应知识点：}\relatedtopic{SC4002-L07-T01}{Head、Dependent 与关系标签}、\relatedtopic{SC4002-L07-T02}{树的约束}。

\begin{workedsolution}
第一句的主要谓词为 \texttt{prefer}，主语是 \texttt{I}，直接宾语是 \texttt{flight}。\texttt{the}、\texttt{morning}、\texttt{Denver} 直接修饰 \texttt{flight}；\texttt{through} 按讲义的 case 关系依存于 \texttt{Denver}。下图是按该页关系重新布局的有向树，不以节点横向位置表示词序。

\begin{center}
% Original layout, relations from Week 7-Dependency Parsing p. 3.
\begin{tikzpicture}[
  word/.style={font=\small\ttfamily,draw=noteBlue!45,rounded corners,fill=noteBlue!4,inner sep=4pt},
  rel/.style={font=\scriptsize,fill=white,inner sep=1.5pt},
  edge/.style={-{Stealth[length=2mm]},thick,draw=noteBlue}]
\node[word] (root) at (0,0) {ROOT};
\node[word] (prefer) at (0,-1.05) {prefer};
\node[word] (i) at (-2,-2.2) {I};
\node[word] (flight) at (1.1,-2.2) {flight};
\node[word] (the) at (-1.2,-3.5) {the};
\node[word] (morning) at (1.1,-3.5) {morning};
\node[word] (denver) at (3.4,-3.5) {Denver};
\node[word] (through) at (3.4,-4.65) {through};
\draw[edge] (root)--(prefer);
\draw[edge] (prefer)--node[rel,above left]{nsubj}(i);
\draw[edge] (prefer)--node[rel,above right]{dobj}(flight);
\draw[edge] (flight)--node[rel,above left]{det}(the);
\draw[edge] (flight)--node[rel,right]{nmod}(morning);
\draw[edge] (flight)--node[rel,above right]{nmod}(denver);
\draw[edge] (denver)--node[rel,right]{case}(through);
\end{tikzpicture}


{\small 原创重绘；依据 SC4002 Week 7 p.~3，标签沿用讲义。}
\end{center}

\texttt{flight} 对 \texttt{prefer} 而言是 dependent，对三个修饰词而言则是 head。\texttt{the} 修饰的是 \texttt{flight}，而不是相邻的 \texttt{morning}；\texttt{through Denver} 的箭头为 \texttt{Denver}$\rightarrow$\texttt{through}，不能按词的前后位置猜方向。

第二句以 \texttt{gives} 为中心：谁给？John；给什么？apple；给谁？Marry。因此
\[
\begin{aligned}
\texttt{gives}&\xrightarrow{\mathrm{nsubj}}\texttt{John}, &
\texttt{gives}&\xrightarrow{\mathrm{obj}}\texttt{apple},\\
\texttt{gives}&\xrightarrow{\mathrm{iobj}}\texttt{Marry}, &
\texttt{apple}&\xrightarrow{\mathrm{det}}\texttt{an}.
\end{aligned}
\]
\texttt{an} 限定 apple，不直接依存于 gives。加入人工节点时另有 ROOT$\rightarrow$gives。主语、直接宾语、间接宾语是相对于动词的不同语法功能，不是三个不同的词性标签。
\end{workedsolution}

\exercisequestion{SC4002-L07-E02}{短语树转换：Vinken will join ...}

\textbf{来源与题意：}讲义 pp.~13--14。根据给定短语树的 head 标记，将 \texttt{Vinken will join the board as a nonexecutive director Nov 29} 转换成词级依存结构。

\textbf{对应知识点：}\relatedtopic{SC4002-L07-T03}{Treebanks 与中心词转换}。

\begin{workedsolution}
先确定局部短语的 lexical heads，再让非中心子短语的 head 依存于当前短语的 head：
\begin{center}
\small
\begin{tabularx}{\textwidth}{@{}l l X@{}}
\toprule
短语或结构 & 给定 head & 对应依存边 \\
\midrule
\texttt{the board} & \texttt{board} & board$\rightarrow$the \\
\texttt{a nonexecutive director} & \texttt{director} & director$\rightarrow$a；director$\rightarrow$nonexecutive \\
\texttt{as a nonexecutive director} & \texttt{director} & director$\rightarrow$as \\
\texttt{Nov 29} & \texttt{29} & 29$\rightarrow$Nov \\
整句及动词短语 & \texttt{join} & join$\rightarrow$Vinken、will、board、director、29 \\
\bottomrule
\end{tabularx}
\end{center}
例如在 $NP\rightarrow DT\ NN$ 中选择 NN：board 代表整个 \texttt{the board}，然后在更高层连接 join$\rightarrow$board。join 作为同一个词级 head 沿 VP、S 向上传递，不生成 join$\rightarrow$join。日期与介词短语的 head 选择按本页原图，不另行改用其他体系；转换需要这些规则，不能只删除短语节点。原树的 NP-SBJ 等功能标记还可帮助赋予依存标签。
\end{workedsolution}

\exercisequestion{SC4002-L07-E03}{Shift-reduce：book me the morning flight}

\textbf{来源与题意：}讲义 p.~17 的完整算法例题。从 $[\mathrm{ROOT}]$ 与全句 buffer 出发，用 Shift、LeftArc、RightArc 构造给定依存树，并说明规约时机和结束条件。

\textbf{对应知识点：}\relatedtopic{SC4002-L07-T04}{Shift-reduce 状态与动作}、\relatedtopic{SC4002-L07-T05}{动作决策}。

\begin{workedsolution}
目标边为
\[
\begin{gathered}
\mathrm{ROOT}\rightarrow\texttt{book},\qquad
\texttt{book}\xrightarrow{\mathrm{iobj}}\texttt{me},\qquad
\texttt{book}\xrightarrow{\mathrm{dobj}}\texttt{flight},\\
\texttt{flight}\xrightarrow{\mathrm{det}}\texttt{the},\qquad
\texttt{flight}\xrightarrow{\mathrm{nmod}}\texttt{morning}.
\end{gathered}
\]
为紧凑展示全部状态，令 $R=\mathrm{ROOT}$、$b=\texttt{book}$、$m=\texttt{me}$、$t=\texttt{the}$、$M=\texttt{morning}$、$f=\texttt{flight}$。\textbf{栈顶在右；每行状态为执行该行动作之后}，不同于讲义原表的操作前状态写法。操作演示暂不预测边标签。

\begin{center}
\small
\begin{tabular}{@{}r l l l l@{}}
\toprule
步骤 & 操作 & Stack & Buffer & 新增关系 \\
\midrule
0 & 初始化 & $[R]$ & $[b,m,t,M,f]$ & --- \\
1 & Shift & $[R,b]$ & $[m,t,M,f]$ & --- \\
2 & Shift & $[R,b,m]$ & $[t,M,f]$ & --- \\
3 & RightArc & $[R,b]$ & $[t,M,f]$ & $b\rightarrow m$ \\
4 & Shift & $[R,b,t]$ & $[M,f]$ & --- \\
5 & Shift & $[R,b,t,M]$ & $[f]$ & --- \\
6 & Shift & $[R,b,t,M,f]$ & $[]$ & --- \\
7 & LeftArc & $[R,b,t,f]$ & $[]$ & $f\rightarrow M$ \\
8 & LeftArc & $[R,b,f]$ & $[]$ & $f\rightarrow t$ \\
9 & RightArc & $[R,b]$ & $[]$ & $b\rightarrow f$ \\
10 & RightArc & $[R]$ & $[]$ & $R\rightarrow b$ \\
\bottomrule
\end{tabular}
\end{center}

\begin{enumerate}
  \item 第 1 步后不能立即 ROOT$\rightarrow$book 并弹出 book：它尚需接收 me 和 flight。先读入 me，建立 book$\rightarrow$me 后可移除 me，因为本例中它没有待接收的子节点。
  \item the 与 morning 都修饰尚未读入的 flight，不能仅因相邻而把两者互连；连续 Shift，直到 flight 入栈。
  \item 栈顶为 morning、flight 时，LeftArc 建立 flight$\rightarrow$morning 并移除 morning；下一次 LeftArc 建立 flight$\rightarrow$the 并移除 the。两次都保留中心词 flight。
  \item flight 的修饰词处理完毕后，RightArc 建立 book$\rightarrow$flight；最后再建立 ROOT$\rightarrow$book。第 6 步 buffer 虽已空，但还需要四次规约。
\end{enumerate}
最终栈只剩 ROOT，buffer 为空，五条目标边全部建立。五个实际词各 Shift 一次、各作为 dependent 移出栈一次，共十次操作；人工 ROOT 不弹出。
\end{workedsolution}
